\chapter{Prefácio}

\fontsize{16}{16}\selectfont

Este livro de exercícios possui como fonte a \textit{Enciclopédia de Tesujis} (手筋大事典), da Nihon Kiin (日本棋院), a Associação Japonesa de Go. Na Coreia do Sul, intitulada como \textit{Sugeundaesajeon} (수근대사전), diz-se que ela foi originalmente de autoria do duo Segoe Kensaku 9p (1889–1972) e um de seus alunos mais famosos, um dos maiores ícones do jogo, Go Seigen 9p (1914-2014). Segoe Kensaku 9p possui seu próprio dicionário de tesujis, porém não está claro se é realmente o caso de ele ter tido alguma influência na enciclopédia atual.

É bem provável que a influência seja concreta no entanto, pois Segoe Kensaku 9p fundou o dojo mais influente de sua época, treinando jogadores que viriam a liderar ou revolucionar o Go, como Hashimoto Utaro 9p, Cho Hunhyun 9p e Go Seigen 9p.

O boca-a-boca sobre a suposta autoria de Segoe é baseada em uma prática que já data de 50 anos pelo menos, no Japão, na Coreia do Sul e na China: sua suposta enciclopédia é uma das principais referências para se treinar aprendizes de profissionais (inseis) até hoje.

\textit{Tesujis} (手筋), ou movimentos habilidosos, são essenciais no Go, e vão além de um brilhantismo de situações específicas. Quando generalizados apropriadamente tornam-se técnicas que ajudam o jogador a evoluir no jogo, não somente a matar ou viver com grupos. Com isso em mente, as edições da Nihon Kiin desta enciclopédia estão capituladas em técnicas e temas. No entanto, não é essa a finalidade desta compilação. Aqui, a organização é um consideravelmente mais aleatória, com o objetivo de instigar ao máximo a leitura do estudante, constituindo, assim, um treinamento de mais alto nível.

A dificuldade dos exercícios nesta coleção abrange uma gama muito grande e bastante mesclada. As diferentes páginas podem ir desde 8 \textit{kyu} até 6 \textit{dan}, mensuradas através da plataforma \href{https://www.101weiqi.com/}{101weiqi}. Com isso em mente, é recomendado não necessariamente querer completar folhas individuais como objetivo de uma sessão de estudos, é talvez mais ideal calibrar o treinamento por tempo ou número de exercícios tentados.
 
\section*{Como Usar Este Livro}

Livros de exercícios não possuem respostas propositalmente. O intuito é desafiar o estudante a estabelecer a melhor variação para ambos os lados.

O estudante resolve o exercício ao escrever números nas intersecções, com os movimentos relevantes. Um professor estrito corrigiria o exercício escrevendo um círculo no topo se a resolução for correta, e uma cruz, se incorreta. Caso o aluno não consiga encontrar os movimentos certos em mais tentativas, o professor pode deixar notas com orientações ou dicas.

\vspace{0.9cm}

\fontsize{14}{14}\selectfont

\hspace{1cm}
\noindent
\begin{minipage}[c]{0.4\linewidth}
  \centering
  \textit{Incorreto}
  \begin{goban}[board dimension       = 12,
                board size            = 19,
                scale                 = 1,
                outline line width    = 0.55mm,
                horizontal clip start = 11,
                horizontal clip end   = 19,
                vertical clip start   = 13,
                vertical clip end     = 19,]
    \parseSgfFile{\repoPath/sgf/sugeundaesajeon/problem-0001_wrong_solution.sgf}
  \end{goban}%
\end{minipage}
\begin{minipage}[c]{0.4\linewidth}
  \centering
  \textit{Correto}
  \begin{goban}[board dimension       = 12,
                board size            = 19,
                scale                 = 1,
                outline line width    = 0.55mm,
                horizontal clip start = 11,
                horizontal clip end   = 19,
                vertical clip start   = 13,
                vertical clip end     = 19,]
    \parseSgfFile{\repoPath/sgf/sugeundaesajeon/problem-0001_correct_solution.sgf}
  \end{goban}%
\end{minipage}

\vspace{1.1cm}

\fontsize{16}{16}\selectfont

Se o leitor não possuir um professor ou jogador mais forte para ajudá-lo, reco-menda-se a utilização da plataforma \href{https://www.101weiqi.com/}{101weiqi}, que possui uma página de procura de problemas (\href{https://www.101weiqi.com/search/}{101weiqi.com/search}) --- às vezes, essa página falha e é preciso procurar o problema através da página de criação de problemas (\href{https://www.101weiqi.com/input/}{101weiqi.com/input}).

Modernamente, livros de problemas de Go padronizam os exercícios para que seja sempre Preto a jogar. Não é o que acontece aqui. Em cada \textit{tsumego}, o leitor deverá avaliar quem jogará primeiro. E os tópicos dos problemas não se restringem somente a vida ou morte, o objetivo é encontrar o melhor resultado para ambos os lados, e isso inclui desde pequenos ganhos de forma a sequências ótimas de fim-de-jogo.

Este não é um livro para iniciantes, e também não para jogadores intermediários. É bastante desafiador mesmo para jogadores avançados, de nível \textit{dan} ou superior. Mas não se intimide pelo desafio, lembre-se de que especialmente errando é que se aprende mais. E especialmente que, estudando, é muito improvável de não caminharmos para frente.

\bigskip
\bigskip

\hspace*{\fill} Philippe Fanaro\hspace{0.02cm}

\hspace*{\fill} Julho de 2026